\section{Derivatives of theta functions and the primes they count}\label{sec:der}

This section is where the primes enter. The values of the composite lift of Theorem~A are Taylor
coefficients (\emph{jets}) of theta functions of the CM curve $\C/\Omega_0\cO_K$ along linear forms
over $E$. At every prime $\ell$ that does not split in $K$ such a jet is divisible by an explicit
power $\ell^{D_\ell}$, the \emph{predicted divisibility}, up to a loss logarithmic in its degree.
For rows whose partners have the orders $\kappa$, and for $\ell\ne d$, $D_\ell$ is the part at
$\ell$ of the prime sum $\sfin$, divided by $2\log\ell$ (\S\ref{der:sec:primesums}).

The mechanism is the following. After a scaling $z\mapsto\varpi z$ the theta function is an
integral series times $\exp(c_\varpi\Xi)$, where $\Xi$ is a combination of \emph{couplings}
$z^{\ell^e}z$. Along a resolution of the identity of $E^n$ the couplings pair each place $\tau$
with its \emph{partner} $-\ell^e\tau$, because for $e\in I_\ell^\circ$ (odd $e$ for inert $\ell$,
all $e$ for $\ell=d$) the $\ell^e$-th power at $\tau$ is, modulo $\pi_\ell$, complex conjugation at
the partner (Lemma~\ref{der:lem:frob}). Killing the couplings by a change of variables costs
determinants, and Legendre's formula gives the rest. (Classically, the curve is supersingular at
such $\ell$.)

From the data of Definition~\ref{setup:def:data} the section obtains four things: the count
(Theorems~\ref{der:thm:count} and~\ref{der:thm:schur}), used in Sections~\ref{sec:ub-divis}
and~\ref{sec:ub-objects}; the integral vacuum, standard data and the archimedean scale
(\S\ref{der:sec:vacuum}), with which Section~\ref{sec:ub-objects} expands the composite lift and
bounds it by $\sinf$; sharpness (Theorem~\ref{der:thm:sharp}), used for the lower bound in
Section~\ref{sec:ol-values}; and the arithmetic of $E$ (\S\ref{der:sec:field}).

\subsection{The field, its places and its primes}\label{der:sec:field}

This subsection fixes the notation for places and for the primes that do not split in $K$. Its
one result, Lemma~\ref{der:lem:frob}, is the congruence that makes the $\ell^e$-th power at a place
agree with complex conjugation at the partner; Definition~\ref{der:def:field} collects arithmetic
data of $E$ that other sections use.

\begin{notation}
Letters are those of Definition~\ref{setup:def:data} and Section~\ref{sec:setup}; in particular
$\Phi\cong\Gal(E/K)\cong(\Z/q_1)^\times$, $M=p_0p_1$ and
$\Omega=e^{\gamma-\Re L'/L(1,\chiK)}2\pi/\rdE$. As there, for an integer $\nu$ and $\tau\in\Phi$,
$\nu\tau$ is the product of $\tau$ with the place of the class of $\nu$, read as $\tau$ if
$\gcd(\nu,q_1)>1$. Each $\tau\in\Phi$ is also an embedding $E\to\C$ of the CM type, and
$\sigma_\tau\in\Gal(E/K)$ is its Galois element; for an integer $\nu$, $\sigma_\nu$ is the Galois
element of the place of the class of $\nu$. For a prime $\ell$ we fix $\iota\colon\C\to\Cl$,
$\iota\mid\ell$; $\cO$ is the ring of integers of $\Cl$, $v$ the valuation with $v(\ell)=1$, and a
series is \emph{integral} if its coefficients lie in $\cO$.
\end{notation}

\begin{definition}[Non-split primes; \lean{New.Parameters.partner}]\label{der:def:nonsplit}
Let $\chiK(\ell)\ne1$, so that $\ell$ is inert or $\ell=d$. Put
$(\pi_\ell,q_\ell,v(\pi_\ell))=(\ell,\ell^2,1)$ if $\ell$ is inert and $(\sqrt{-d},\ell,\frac12)$ if
$\ell=d$, and put $\rho_\ell=v(\pi_\ell)/(q_\ell-1)$ and $\gamma_\ell=\frac1{\ell-1}-\rho_\ell$. Let
$I_\ell^\circ$ be the set of odd numbers if $\ell$ is inert and $\N$ if $\ell=d$, and
$I_\ell=I_\ell^\circ\setminus\{0\}$; then $\gamma_\ell=v(\pi_\ell)\sum_{e\in I_\ell}\ell^{-e}$. The
\emph{partner} of $\tau\in\Phi$ for $e$ is $-\ell^e\tau$. For split $\ell$, $I_\ell=\emptyset$ and
$\gamma_\ell=0$. As primes dividing $q_1$ split (\lean{New.Parameters.chiK_eq_one_of_dvd}), a
non-split $\ell$ does not divide $q_1$.
\end{definition}

\begin{lemma}[Frobenius; \lean{New.Parameters.norm_pow_sub_conj_partner_le}, \lean{New.Parameters.norm_conj_pow_sub_partner_le}]
\label{der:lem:frob}
Let $\chiK(\ell)\ne1$, $e\in I_\ell^\circ$, $\tau\in\Phi$, and let $\tau'=-\ell^e\tau$ be its partner.
Then in $\Cl$ we have $\tau(w)^{\ell^e}\equiv\overline{\tau'(w)}\pmod{\pi_\ell}$ for $w\in\cO_E$, and
$\overline{\tau(w)}^{\,\ell^e}\equiv\tau'(w)$ for $w\in E$ integral above $\ell$.
\end{lemma}

\begin{proof}
The maps $w\mapsto\tau(w)^{\ell^e}$ and $w\mapsto\tau(\sigma_{-\ell^e}\bar w)$ from
$\cO_E=\Z[\beta_K,\zeta]$, $\beta_K=\frac{1+\sqrt{-d}}2$, to $\cO/\pi_\ell$ are ring maps, so it
suffices that they agree on $\zeta$, where $\sigma_{-\ell^e}(\bar\zeta)=\zeta^{\ell^e}$, and on
$\sqrt{-d}$ (\lean{New.Parameters.forall_map_galOfPlaceInt_eq_pow_iff}). For inert $\ell$,
$(\sqrt{-d})^{\ell^e}\equiv\chiK(\ell)^e\sqrt{-d}=-\sqrt{-d}$ as $e$ is odd; for $\ell=d$ both sides
are $\equiv0\bmod\sqrt{-d}$. For $w$ integral above $\ell$ write $w=w_0/m'$ with $w_0\in\cO_E$ and
$\ell\nmid m'$, and use ${m'}^{\ell^e}\equiv m'$.
\end{proof}

\begin{definition}[Arithmetic data used elsewhere]\label{der:def:field}
(a) \emph{Chosen places} (\lean{New.Data.IsChosen}, \lean{New.Data.ChosenPlace}): let
$\bar\pi_M\in\cO_K$ have norm $M$ and let $\pi_M$ (\lean{New.Data.piM}) be the conjugate of its
image in $\cO_E$, so that $\pi_M\bar\pi_M=M$ (\lean{New.Data.piM_mul_star}). The chosen places are
the primes dividing $\pi_M$; each has norm $p_0$ or $p_1$ (\lean{New.Data.q_eq_levelPrime}), and
those above $p_i$ correspond to $\Phi$ (\lean{New.Parameters.nonempty_infinitePlace_equiv_finPlace}).
(b) \emph{Characters}: $\omega_{\fq}$ (\lean{New.Data.chiMLocal}) is the $q_2$-th power of a fixed
character of order $q_2^2$ of the residue field at a chosen $\fq$, and is trivial elsewhere;
$\mu^\flat$ (\lean{New.Data.chiZero}) and $\chi$ (\lean{New.Data.chiM}) are the automorphic characters
of $\U(1)$ of Definition~\ref{data:def:liftchar}, the first of weight $-(\kappa+\Delta/2)$ and trivial
tame part, the second of weight $0$ and tame part $\omega$.
(c) \emph{Critical pairs} (\lean{New.Data.CriticalPair}): a character $\nu$ of weight $-c$ whose
tame part is supported on the chosen places, and an integer $z\in[1,t]$ with $z\le c(\tau)\le4b$
at every $\tau$; there are finitely many. The letters $\nu$, $z$, $c$ are those of
Section~\ref{sec:det-chain}, which defines ramified pairs, main pairs and their ratios
(\S\S\ref{det-chain:sec:setup} and~\ref{det-chain:sec:regular}).
(d) \emph{The different}: $\delta_E=u\,P'(\zeta)\sqrt{-d}$ (\lean{New.Parameters.deltaE}), with
$P'$ the derivative of the $q_1$-th cyclotomic polynomial and $u\bar u=1$, is a purely imaginary
generator of the different of $E$. It has $\prod_\tau|\tau(\delta_E)|=\rdE^f$
(\lean{New.Parameters.prod_deltaE_eq_rootDiscr_pow}), and $|\tau(\delta_E)|,\rdE\in
[x^{10},x^{12}]$ (\lean{New.Data.hasLogBound_deltaE}, \lean{New.Data.rootDisc_mem_Icc}).
(e) \emph{Weil data} (\lean{New.Parameters.instWeilDataE}): $\jmath=-1/\sqrt{-d}$, the idele
$\alpha$ with component $1/P'(\zeta)$ at the places dividing $\beta_K$ and $1$ elsewhere, and a
splitting character $\upsilon$. Put $\fa_0=E\cap\alpha\hat\cO$; then $\jmath\fa_0\bar\fa_0$ is the inverse
different. The number $s_0$ (\lean{New.Parameters.vacuumRoot}), an $f$-th root of
$\upsilon(\alpha)^{-1}\abs{\alpha}^{1/2}$, $\abs{\alpha}$ the idele norm, has $|\sigma s_0|^2=\rdE/\sqrt d$ for $\sigma\in\Gal(\C/K)$
(\lean{New.Parameters.norm_gal_vacuumRoot}).
\end{definition}

\subsection{The curve and its exponent at a non-split prime}\label{der:sec:curve}

This subsection works on a single CM curve. It recalls the formal group of the curve at a
non-split prime (Fact~\ref{der:fact:formal}) and proves Theorem~\ref{der:thm:exponent}: after the
scaling $z\mapsto\varpi z$ the theta function is an integral series times $\exp(c_\varpi\Xi)$, with
$\Xi$ congruent modulo $\pi_\ell$ to a combination of the $\ell^j$-th powers of the couplings
$z^{\ell^e}z$. From it Proposition~\ref{der:prop:thetaprod} passes to products of theta functions
in many variables (\S\ref{der:sec:normal}).

A \emph{CM lattice} $C$ is a lattice $\Lambda\subset\C$ stable under $\cO_K$ with a point $t_0$ of
order two and $\wp''(t_0)=2$. Its theta function $\theta$ is the translate by $t_0$ of the sigma
function, twisted to a theta function, with $\theta(0)=1$. The functions
$\mathrm x=\wp(z+t_0)-\wp(t_0)=z^2+\cdots$ and $\mathrm y=\frac12(\wp'(z+t_0)-a_1\mathrm x)$
satisfy $\mathrm y^2+a_1\mathrm x\mathrm y=\mathrm x^3+\mathrm x$, and we put $g=-\log\theta$.
Taylor series at $0$ are read in $\Cl$ through $\iota$, and $\cO[[\mathrm y]]$ is the ring of
integral series in $\mathrm y$. The curve of the proof is $\C/\Omega_0\cO_K$
(\lean{New.Parameters.cmCurve}).

\begin{fact}[Formal group; \lean{New.CMLattice.taylorVia_translatedX_mem_integralSeriesRing},
\lean{New.CMLattice.derivative_mem_integralSeriesRing},
\lean{New.CMLattice.mem_integralSeriesRing_iff_exists_subst_formalParam},
\lean{New.CMLattice.exists_scaled_formalParam}]\label{der:fact:formal}
Let $d\equiv7\bmod8$ be prime. (a) The ring $\cO[[\mathrm y]]$ contains $\mathrm x$ and is stable
under $d/dz$. (b) The part $\mathcal U$ of $\mathrm y$ in the degrees $\equiv1\bmod q_\ell-1$ has
$\cO[[\mathcal U]]=\cO[[\mathrm y]]$. (c) For some $\varpi$ with $v(\varpi)=\rho_\ell$, the series
$\mathcal Y=\mathcal U(\varpi z)/\varpi$ is integral and
$\mathcal Y\equiv\sum_{j\ge0}z^{q_\ell^j}\pmod{\pi_\ell}$.
(\leanfile{Lib/CM/FormalGroup/IntegralParameter}, \leanfile{Lib/CM/FormalGroup/FormalParameter};
after Lubin--Tate, Formal complex multiplication in local fields.)
\end{fact}

Hence $h(\varpi z)$ is integral and $v(h^{(n)}(0))\ge v(n!)-n\rho_\ell\ge\gamma_\ell n-O(\log n)$
for $h\in\cO[[\mathrm y]]$; this is the rate of Katz (Divisibilities, congruences, and Cartier
duality, 1981). Put $c_\varpi=\varpi^{\ell+1}/\ell$. Then $v(c_\varpi)<0$ but
$v(c_\varpi\pi_\ell)=\frac1{\ell-1}$ (\lean{New.Parameters.norm_pow_succ_div_mul_piPrimeInt}), so
$\exp(c_\varpi\pi_\ell G)$ is integral for integral $G$ with $G(0)=0$.

\begin{theorem}[The exponent on one curve; \lean{New.CMLattice.exists_exponent_tsum}]
\label{der:thm:exponent}
Let $\chiK(\ell)\ne1$. There are $\varpi$ with $v(\varpi)=\rho_\ell$, integral $\mathcal W,\Xi$
with $\Xi(0)=0$ and $\lambda_{ej}\in\cO$ ($e\in I_\ell^\circ$, $j\ge0$), the $\lambda_{e0}$ units if
$\ell$ is inert, with
\[
\theta(\varpi z)=\mathcal W(z)\exp\bigl(c_\varpi\Xi(z)\bigr),\qquad
\Xi\equiv\textstyle\sum_{e\in I_\ell^\circ}\sum_{j\ge0}\lambda_{ej}\,(z^{\ell^e}z)^{\ell^j}\pmod{\pi_\ell}.
\]
\end{theorem}

\begin{proof}
\emph{Frobenius congruence} (\lean{New.CMLattice.exists_rescale_theta_eq_pow_mul_one_add_smul}):
$\theta(\pi_\ell z)=\theta(z)^{q_\ell}(1+\pi_\ell H)$ with $H\in\cO[[\mathrm y]]$. As $\pi_\ell$ is an
isogeny of degree $q_\ell$, the quotient is an even elliptic function with integral torsion values,
a polynomial $R(\mathrm x)$, $2\deg R<q_\ell$. It is $\exp(q_\ell g(z)-g(\pi_\ell z))$, whose
coefficients $(q_\ell-\pi_\ell^n)g_n$ are divisible by $\pi_\ell$ for $n<q_\ell$, since
$\pi_\ell^2\mid q_\ell-\pi_\ell^n$ ($n\ge2$) and $\pi_\ell g_n$ is integral
($g''=g''(0)+\mathrm x$, $v(n(n-1))\le v(\pi_\ell)$).
So $R\equiv1\bmod\pi_\ell$, and $q_\ell g(z)-g(\pi_\ell z)\in\pi_\ell\cO[[\mathrm y]]$.
\emph{Averaging}: for the part $g^\flat$ of $g$ in the degrees $\equiv\ell+1\bmod q_\ell-1$,
$\exp(g^\flat-g)\in\cO[[\mathrm y]]$, as $\ell\nmid q_\ell-1$ and its $(q_\ell-1)$-th power is the
product of the $\exp(\zeta'^{-\ell-1}g(\zeta'z)-g(z))$, $\zeta'^{q_\ell-1}=1$. These lie in the
closed ring $\cO[[\mathrm y]]$: $\exp(\alpha\bar\alpha g-g(\alpha z))=\theta(\alpha z)/\theta^{\alpha\bar\alpha}$
does for $\alpha\in\cO_K$, and
$(\zeta'^{-1},\zeta'^{-\ell-1})=\lim(\alpha_0^{q_\ell^n},(\alpha_0\bar\alpha_0)^{q_\ell^n})$ for
$\alpha_0\in\cO_K$ lifting $\zeta'^{-1}$, as $\bar\alpha_0\equiv\alpha_0^\ell\bmod\pi_\ell$; then
$z\mapsto\zeta'z$
(\lean{New.mem_normPairClosure_of_pow_eq},
\lean{New.CMLattice.expSubst_rescale_sub_mem_of_pow_eq_one}). So
$\theta(\varpi z)=\mathcal W\exp(c_\varpi\Xi_1)$ with $\Xi_1=-g^\flat(\varpi z)/c_\varpi$.
\emph{Functional equation} (\lean{New.CMLattice.exists_functional_equation_exponentSeries}): the
congruence for $g^\flat$ at $\varpi z$ gives
$(1-\pi_\ell^n/q_\ell)[z^n]\Xi_1\equiv\lambda[z^n]\mathcal Y^{i_0}$, $\lambda\in\cO$, with
$i_0=\ell+1$ (inert) or $2$ ($\ell=d$). As $\Xi_1$ has no terms below $i_0$ and
$\pi_\ell^n/q_\ell\equiv0$ for $n>2$, $\Xi_1\equiv\lambda\mathcal Y^{\ell+1}$ (inert),
$\Xi_1\equiv\lambda(\mathcal Y^2-z^2/2)$ ($\ell=d$, where $1-\pi_\ell^2/q_\ell=2$).
\emph{Solution}: by Fact~\ref{der:fact:formal}(c), for inert $\ell$,
$\mathcal Y^\ell\mathcal Y\equiv\sum_{i\text{ odd}}z^{\ell^i}\sum_{i\text{ even}}z^{\ell^i}=\sum_{e\text{
odd}}\sum_j(z^{\ell^e}z)^{\ell^j}$, and for $\ell=d$, $\mathcal Y^2\equiv\sum_j(z^2)^{\ell^j}+2\sum_{e\ge1,j}
(z^{\ell^e}z)^{\ell^j}$. For inert $\ell$, $\lambda$ is a unit multiple of $h_1=[z^{\ell-1}]\mathrm x$, and
$h_n=[z^{\ell-1}]\mathrm x^n$ satisfies $(4n+6)h_{n+2}+(n+1)a_1^2h_{n+1}+(4n+2)h_n\equiv0$ (the
$z^{\ell-1}$-coefficient of $(\mathrm x^n\mathrm x')'$), which leads from $h_1$ to $h_{(\ell-1)/2}=1$; so
$h_1$ is a unit (\lean{New.CMLattice.norm_coeff_taylorVia_translatedX_eq_one}).
\end{proof}

\subsection{The normal form and products of theta functions}\label{der:sec:normal}

This subsection passes from one variable to many. It defines couplings and the normal form, shows
that a series in normal form becomes integral in low degrees after a substitution that makes the
couplings divisible (Lemma~\ref{der:lem:nf}), and proves that products of theta functions along a resolution
of the identity by vectors integral above $\ell$ have the normal form
(Proposition~\ref{der:prop:thetaprod}). These two properties
are what the count of \S\ref{der:sec:count} uses.

For variables $z_{\tau,i}$ ($\tau\in\Phi$, $i$ in a finite set), a non-split $\ell$ and
$e\in I_\ell^\circ$, the \emph{coupling} is
$\Gamma_e=\sum_{\tau,i}z_{\tau,i}^{\ell^e}z_{-\ell^e\tau,i}$ (\lean{New.coupling}); it is
homogeneous of degree $\ell^e+1$.

\begin{definition}[Normal form; \lean{New.HasNormalForm}]
A series $F$ has the \emph{normal form} at $\ell$ (for $\varpi$ and $\lambda_{ej}\in\cO$) if
$F(\varpi z)=\mathcal W\exp(c_\varpi\Xi)$ with $\mathcal W$ integral and
$\Xi=\sum_{e\in I_\ell^\circ}\sum_j\lambda_{ej}\Gamma_e^{\ell^j}$.
\end{definition}

\begin{lemma}[\lean{New.HasNormalForm.norm_coeff_subst_le_one}, \lean{New.HasNormalForm.of_dvd_sub}]
\label{der:lem:nf}
Let $F$ have the normal form, let $Z$ be integral without constant term, and let $\pi\in\cO$ with
$v(c_\varpi\pi)\ge\frac1{\ell-1}$ divide $\Gamma_e(Z)$ for $e\in I_\ell^\circ$ with $\ell^e\le B$.
Then $F(\varpi Z)$ is integral in the degrees $\le B$ (the terms of $\Xi(Z)$ with $\ell^e>B$ have
order $>B$, and the others are divisible by $\pi$). For the normal form it suffices that
$F(\varpi z)=\mathcal W\exp(c_\varpi\Xi')$ with $\mathcal W,\Xi'$ integral, $\Xi'(0)=0$ and
$\Xi'\equiv\Xi\bmod\pi$.
\end{lemma}

A \emph{resolution of the identity} (\lean{New.IsResolution}) is a family of vectors $u_j\in E^n$
with multiplicities $c_j$ such that $\sum_jc_j(u_j^\dagger z)u_j=z$ on $\C^{\Phi\times n}$, where
$u^\dagger z=\sum_{\tau,i}\overline{\tau(u_i)}z_{\tau,i}$. A \emph{tight frame} of length $M'$
(\lean{New.IsTightFrame}) is a family $w_i\in\fa_0$ with multiplicities $c_i$ and
$\sum_ic_i\tau(w_i)\overline{\sigma(w_i)}=M'[\tau=\sigma]$. If $M'=b'^2$ and $c_i=1$, the vectors
$w_i/b'$, placed in each coordinate, form a resolution (\lean{New.IsTightFrame.isResolution}).

\begin{proposition}[Normal form of theta products; \lean{New.Parameters.exists_hasNormalForm}]
\label{der:prop:thetaprod}
Let $\chiK(\ell)\ne1$. There are $\varpi$, $v(\varpi)=\rho_\ell$, and $\lambda_{ej}\in\cO$ (the
$\lambda_{e0}$ units if $\ell$ is inert) such that for every resolution $(u_j,c_j)$ by vectors
integral above $\ell$ and every integral series $H$ in the $\mathrm y(u_j^\dagger z)$, the product
$\prod_j\theta(u_j^\dagger z)^{c_j}\cdot H$ has the normal form.
\end{proposition}

\begin{proof}
Let $L_j=u_j^\dagger z$, a form with integral coefficients. By Theorem~\ref{der:thm:exponent},
$\prod_j\theta(\varpi L_j)^{c_j}=\prod_j\mathcal W(L_j)^{c_j}\exp(c_\varpi\sum_jc_j\Xi(L_j))$. By
Lemma~\ref{der:lem:frob} and additivity of $\ell^e$-th powers mod $\ell$,
$L_j^{\ell^e}\equiv\sum_{\tau,i}\tau'(u_{j,i})z_{\tau,i}^{\ell^e}$ with $\tau'=-\ell^e\tau$, so by the
resolution identity (\lean{New.Parameters.piPrimeInt_dvd_sum_linearSeries_sub_coupling})
\[
\textstyle\sum_jc_jL_j^{\ell^e}L_j\equiv\sum_{\tau,i,\sigma,i'}z_{\tau,i}^{\ell^e}z_{\sigma,i'}
\sum_jc_j\tau'(u_{j,i})\overline{\sigma(u_{j,i'})}=\Gamma_e\pmod{\pi_\ell}.
\]
As $\ell^{j'}$-th powers are additive mod $\pi_\ell$,
$\sum_jc_j(L_j^{\ell^e}L_j)^{\ell^{j'}}\equiv\Gamma_e^{\ell^{j'}}$, so
$\sum_jc_j\Xi(L_j)\equiv\sum_{e,j'}\lambda_{ej'}\Gamma_e^{\ell^{j'}}$; apply Lemma~\ref{der:lem:nf}, as
$H(\varpi z)$ is integral.
\end{proof}

\subsection{The integral vacuum and standard data}\label{der:sec:vacuum}

This subsection connects the theta functions of the curve with the vectors of the Weil
representation from which the lifts are built. It defines the jet scale, the integral vacuum and
standard data (Definition~\ref{der:def:std}); it shows that the integral vacuum has integral values
with bounded conjugates (Proposition~\ref{der:prop:vacuum}) and that the jet series of standard data
are theta products (Proposition~\ref{der:prop:product}), whence a lower bound for the valuation of
their jets and, at a non-split prime and for level one, the normal form
(Corollary~\ref{der:cor:integral}); and it
computes the absolute value of the jet scale and describes the action of $\Gal(\C/K)$ on jets
(Proposition~\ref{der:prop:scale}). The count uses Corollary~\ref{der:cor:integral}(b);
Section~\ref{sec:ub-objects} uses Corollary~\ref{der:cor:integral}(a), standard data and
Proposition~\ref{der:prop:scale}.

For a hermitian space $V$ of rank $n$ over $E$ the background provides the space $\mathscr S(V)$
of the adelic Weil representation, the theta functional $\theta_V$, the translations
$\mathrm H(\mathbf y)$ and the \emph{vacuum} $\mathrm{vac}_V$ (Gaussian at infinity times twisted
indicator of the standard lattice). A vector $\mathcal F\in\mathscr S(V)$ is read as the function
$\mathbf y\mapsto\theta_V(\mathrm H(\mathbf y)\mathcal F)$. The space $V$ is \emph{standard} if its
signs are $+1$ and its point is $1$. Put $\mathrm{finVec}(\mathbf x)=(0,\alpha\mathbf x)$; the finite
Heisenberg operator $\mathrm H_{\mathbf x}$ (\lean{New.Space.finHeisenbergOp}) is
$\mathrm H(\mathrm{finVec}\,\mathbf x)$ times a root of unity, and the \emph{finite part}
$\mathscr S_{\mathrm f}(V)$ is the span of the $\mathrm H_{\mathbf x}\mathrm{vac}_V$. The \emph{jet
series} $J_\eta\mathcal F$ (\lean{New.Space.jetSeries}), $\eta\in\C$, has coefficient
$\eta^{|\mathbf k|}\theta_V(\mathsf v_{z^{\mathbf k}}\otimes\mathcal F)/\mathbf k!$ at $z^{\mathbf k}$,
where $\mathsf v_{z^{\mathbf k}}$ is the Fock vector of $z^{\mathbf k}$. On the finite part it is the
Taylor series at $0$ of $z\mapsto\mathcal E_{\mathcal F}(\eta z)$, where $\mathcal E_{\mathcal F}$ is
the entire extension of $\mathbf y\mapsto\theta_V(\mathrm H(\mathbf y)\mathcal F)e^{|\mathbf y|^2/2}$.
For $P$ homogeneous of degree $l$ we write $\mathrm{jet}_P(J_\eta\mathcal F):=(P(\partial)J_\eta\mathcal F)(0)$;
it equals $\eta^l\theta_V(\mathsf v_P\otimes\mathcal F)$ (\lean{New.Space.jet_jetSeries_of_isHomogeneous}).
For $\eta\ne0$, $J_\eta$ is injective on the finite part (\lean{New.Space.eq_zero_of_jetSeries_eq_zero}).

\begin{definition}[Scales, integral vacuum, standard data]\label{der:def:std}
The \emph{jet scale} (\lean{New.Parameters.jetScale}) is
$\varsigma=s_0u_{\mathrm{CS}}^{-1}\sqrt{2\pi/\sqrt d}/\Omega_0$, where $u_{\mathrm{CS}}$ is the unit
of Fact~\ref{der:fact:bg}(a) for $\C/\Omega_0\cO_K$. The \emph{vacuum scale} $c_s\ne0$
(\lean{New.Parameters.vacuumScale}) makes the values of $c_s\mathrm{vac}$ in rank one at the points
$\mathrm{finVec}(\mathbf x)$ integral at every $\iota$, with conjugates over $K$ of modulus
$\le e^{C_{\mathrm{vac}}q_1^2}$ ($C_{\mathrm{vac}}$ = \lean{New.Parameters.vacuumBoundConst}). The
\emph{integral vacuum} is $\theta^\circ_V=c_s^n\mathrm{vac}_V$ (\lean{New.Parameters.intVacuum}). For
$\delta\in\cO_K$ the \emph{standard data} of level $\delta$ at $\ell$ (\lean{New.Parameters.standardData})
are the span over $\cO_\iota$ (the numbers integral at $\iota$, Section~\ref{sec:setup}) of the
$\mathrm H_{\mathbf x}\theta^\circ_V$ with $\delta\mathbf x$ integral above $\ell$
(\lean{New.Parameters.integralAfter}).
\end{definition}

\begin{fact}[\lean{New.CMLattice.chowla_selberg}, \lean{New.Space.exists_entireTheta_eq_of_isThetaFunction},
\lean{New.CMLattice.exists_rescale_twistedTranslate_eq_theta_pow_mul}]\label{der:fact:bg}
(a) For a CM lattice $C$ there is $u_{\mathrm{CS}}\in\C$, a unit at every $\iota$, with
$\operatorname{area}(C^\sigma)|\sigma u_{\mathrm{CS}}|^2=\pi e^{\gamma-\Re L'/L(1,\chiK)}2\pi/\sqrt d$
for $\sigma\in\Gal(\C/K)$ (Selberg--Chowla, On Epstein's zeta-function, 1967). (b) For $V$ standard
and integral $\delta\in E$ with $\delta\bar\delta\in\N\setminus\{0\}$, every entire function fixed by
the signed translations along $\bar\delta\fa_0^n$ is $\mathcal E_{\mathcal F}$ for a vector
$\mathcal F$ of the finite part (Mumford, Tata lectures on theta I). (c) If $\bar\delta$ is an
isogeny of degree $N_\delta$ of $C$ and $\delta\eta\in\Lambda$, the series of
$(T_\eta\theta)(\bar\delta z)$, $T_\eta$ the twisted translation, is $\theta^{N_\delta}$ times an
element of $\cO[[\mathrm y]]$ (Silverman, The Arithmetic of Elliptic Curves, ch.~IV and VII).
\end{fact}

\begin{proposition}[\lean{New.Parameters.norm_intVacuum_le_one}, \lean{New.Parameters.norm_ringEquiv_intVacuum_le},
\lean{New.Parameters.exists_intVacuum_ne_zero},
\lean{New.Parameters.norm_dataInner_self_le}, \lean{New.Parameters.finHeisenbergOp_mem_standardData},
\lean{New.Parameters.omega_mem_standardData}]\label{der:prop:vacuum}
Let $V$ be standard. The values $\theta^\circ_V(\mathrm{finVec}\,\mathbf x)$ are integral at every
$\iota$, not all zero, with all conjugates under $\Aut(\C)$ of modulus $\le e^{nC_{\mathrm{vac}}
q_1^2}$; if $|\mathcal F(\mathrm{finVec}\,\mathbf x)|\le A$ for all $\mathbf x$, then
$\langle\mathcal F,\mathcal F\rangle\le A^2$. Standard data of level $\delta$ are stable under
$\mathrm H_{\mathbf x}$ for $\delta\mathbf x$ integral above $\ell$, and under $\omega(g)$ for unitary
$g$ integral above $\ell$.
\end{proposition}

\begin{proof}
Values in rank $n$ are products of values in rank one. There (Fact~\ref{der:fact:bg}(b)) a
multiple of the vacuum is a sum of $q_1^{2f}$ translates of a vector whose values are products of
torsion values of $\theta$, integral with bounded conjugates
(\lean{New.Parameters.exists_mul_vacElem_integral_bounded}). Finally, $\Aut(\C)$ acts on absolute
values through $\Gal(\C/K)$ (\lean{New.Parameters.exists_algEquiv_norm_eq},
\lean{New.Parameters.norm_gal_intVacuum_le}).
\end{proof}

\begin{proposition}[Standard data are theta products;
\lean{New.Parameters.exists_jetSeries_map_eq_thetaProd_pow_mul}]\label{der:prop:product}
Let $V$ be standard, $0\ne\delta\in\cO_K$, $N_\delta=\delta\bar\delta$, and let $\mathcal F$ be
standard data of level $\delta$ at $\ell$. There are a resolution $(u_j,c_j)$ of $E^n$ with
$\alpha^{-1}u_j$ integral above $\ell$ and an integral series $H$ in the $\mathrm y(u_j^\dagger z)$
with $\iota J_{\bar\delta\sqrt{2\pi/\sqrt d}/\Omega_0}\mathcal F=\bigl(\prod_j\theta(u_j^\dagger
z)^{c_j}\bigr)^{N_\delta}H$.
\end{proposition}

\begin{proof}
(1) \emph{Vacuum expansion} (\lean{New.Parameters.exists_isVacuumExpansion}): for a tight frame
$(w_i)$ of length $b'^2$, $\ell\nmid b'$, the product of the $\theta(u^\dagger z)$, $u=w_i/b'$ in
each coordinate, is by the frame identity a theta function for $b'\fa_0^n$, so it is
$\mathcal E_{\mathcal F_0}$ with $\mathcal F_0$ in the finite part (Fact~\ref{der:fact:bg}(b)).
Rational translates of $\mathcal F_0$ give products of twisted translates of $\theta$, and Fourier
inversion gives $\theta^\circ_V=\sum_j\gamma_j\mathrm H(\mathrm{finVec}\,\mathbf x_j)\mathcal F_0$ with
$\gamma_j\in\cO_\iota$ and $b'\mathbf x_j$ integral. (2) A translate of $\mathcal F_0$ by $\mathbf x$
with $\delta\mathbf x$ integral above $\ell$ is one by a rational $\mathbf r$ with
$m'\delta\mathbf r\in\fa_0^n$, $\ell\nmid m'$, so its series is a product of
$(T_\eta\theta)(\bar\delta\,\cdot)$ with $m'b'\delta\eta$ periods. (3) Apply
Fact~\ref{der:fact:bg}(c), and take as resolution the $w_i/(b'N')$ with multiplicity ${N'}^2$, for
$N'$ prime to $\ell$ a common multiple of the $m'b'$ over the generators of $\mathcal F$.
\end{proof}

\begin{corollary}[\lean{New.Parameters.norm_jet_jetSeries_le_pow}, \lean{New.Parameters.norm_jet_jetSeries_le_one},
\lean{New.Parameters.exists_hasNormalForm_jetSeries}]
\label{der:cor:integral}
Let $V$ be standard and $\mathcal F$ standard data of level $\delta\ne0$ at $\ell$. (a) For
$\eta\in\C$ and $P$ homogeneous of degree $l$ over $\cO_\iota$,
$v(\iota\,\mathrm{jet}_P(J_{\eta\varsigma}\mathcal F))\ge l\,v(\iota(\eta/(\bar\delta d)))$. (b) If
$\chiK(\ell)\ne1$, $\delta=1$ and $v(\iota\eta)\ge0$, then $\iota J_{\eta\varsigma}\mathcal F$ has the
normal form at $\ell$ for some $\varpi$ with $v(\varpi)=\rho_\ell$.
\end{corollary}

\begin{proof}
(a) As $s_0u_{\mathrm{CS}}^{-1}\overline{\tau(w)}$ is integral when $\alpha^{-1}w$ is integral above
$\ell$ (\lean{New.Parameters.norm_periodFactor_mul_daggerForms_le_one}),
Proposition~\ref{der:prop:product} puts $J_{\bar\delta d\varsigma}\mathcal F$ in the ring generated
by the integral series in the $\mathrm y(\mu_jz)$, the $\theta(\mu_jz)$ and the $g'(\mu_jz)$, where
$\mu_j=d\nu_j$ with $\nu_j$ integral. This ring has integral constant terms and is stable under
$\partial$, as $d\,g''\in\cO[[\mathrm y]]$; so its jets are integral
(\lean{New.CMLattice.norm_jet_thetaProd_pow_mul_le_one}).
(b) As $\ell\nmid q_1$, $s_0$ is an $\ell$-adic unit and the $u_j$ are integral above $\ell$
(\lean{New.Parameters.isIntegralAbove_of_finiteImage_mem}); apply Proposition~\ref{der:prop:thetaprod}
and rescale by $\eta$ (\lean{New.HasNormalForm.rescale}).
\end{proof}

The Galois part of the next proposition rests on the main theorem of complex multiplication, in
the following form.

\begin{fact}[Complex multiplication for theta functions; \lean{New.CMLattice.complex_multiplication_theta}]
\label{der:fact:cm}
Let $C$ be a CM lattice whose lattice is $\Omega_0\cO_K$, and let $\sigma\in\Gal(\C/K)$. There are
$\Omega_\sigma\in\C^\times$, a non-zero integral ideal $\mathfrak c$ of $\cO_K$ and a finite adele
$\beta_\sigma$ of $K$ such that the lattice of the conjugate curve $C^\sigma$ is
$\Omega_\sigma\mathfrak c^2/\Nm\mathfrak c$, has area $|\Omega_\sigma|^2\sqrt d/2$ and contains a
positive integral multiple of $\Omega_\sigma$, and such that for each $\phi\in K$ and each $\phi'\in K$
close enough to $\beta_\sigma\phi$ in the finite adeles, $\sigma$ takes the Taylor series of the
twisted translate $T_{\phi\Omega_0}\theta$ to that of the twisted translate
$T_{\phi'\Omega_\sigma}\theta^\sigma$ of the theta function $\theta^\sigma$ of $C^\sigma$. This is a
case of the main theorem of complex multiplication without the reciprocity law: $\beta_\sigma$ is
not identified with the Artin map and is not shown to be a unit (\leanfile{Lib/CM/ChowlaSelberg};
after Shimura, Introduction to the Arithmetic Theory of Automorphic Functions, 1971, Thm.~5.4, and
Milne, Complex Multiplication, 2020, 9.17).
\end{fact}

\begin{proposition}[\lean{New.Parameters.norm_jetScale_sq}, \lean{New.Data.norm_conj_piMBar_mul_jetScale_sq},
\lean{New.Parameters.exists_isJetConjugate_translates}]\label{der:prop:scale}
We have $|\varsigma|^2=\Omega^{-1}$ and $|\pi_M\varsigma|^2=M/\Omega$. Every standard $V$ has
$\mathcal F$ in the finite part, $\theta_V(\mathcal F)=1$, whose translates span it, such that each
$\sigma\in\Gal(\C/K)$ carries the jets of $\mathrm H(\mathrm{finVec}\,\mathbf x)\mathcal F$ at the
scale $\varsigma$ to those of $\mathrm H(\mathrm{finVec}\,\mathsf B_\sigma\mathbf x)\mathcal F'$ at a
scale $\varsigma'$ with $|\varsigma'|^2=\Omega^{-1}$ (places permuted by $\sigma$), for some
$\mathcal F'$, $\varsigma'$ and a finite adele $\mathsf B_\sigma$ independent of $V$.
\end{proposition}

\begin{proof}
From $|s_0|^2=\rdE/\sqrt d$, $\operatorname{area}(\Omega_0\cO_K)=|\Omega_0|^2\sqrt d/2$ and
Fact~\ref{der:fact:bg}(a) one computes $|\varsigma|^2=\rdE/(2\pi e^{\gamma-\Re L'/L(1,\chiK)})$.
For the Galois part, $\mathcal F$ is the product of $\theta$ along $\zeta^i/q_1$. The series of a
rational translate of $\mathcal F$ is a product of translates of $\theta$ by points $\phi\Omega_0$,
$\phi\in K$. By Fact~\ref{der:fact:cm}, $\sigma$ takes these to the translates by $\phi'\Omega_\sigma$,
$\phi'\in K$ near $\beta_\sigma\phi$ ($\beta_\sigma$ a finite adele), of the $\theta$ of the conjugate
curve, of area $|\Omega_\sigma|^2\sqrt d/2$, at a scale $\varsigma'$ with $|\varsigma'|^2=\Omega^{-1}$
as above. Local constancy and density of rational points give all $\mathbf x$, with $\mathsf B_\sigma$
the image of $\beta_\sigma$.
\end{proof}

So a jet of order $j$ at the scale $\pi_M\varsigma$ has at every conjugate a factor $(M/\Omega)^{j/2}$,
whence the term $j\log(M/\Omega)$ of $\sinf(j)=\log j!+j\log(M/\Omega)$
(Definition~\ref{setup:def:primesums}, \lean{New.Data.jetSize}).

\subsection{The count}\label{der:sec:count}

This subsection proves the main estimate of the section. For a series in normal form at a
non-split $\ell$, with $v(\varpi)=\rho_\ell$, the jets
along products of minors have valuation at least the predicted divisibility $D_\ell$ of the
diagram, less a loss logarithmic in the number of boxes (Theorem~\ref{der:thm:count}).
Theorem~\ref{der:thm:schur} is the form in which Sections~\ref{sec:ub-divis}
and~\ref{sec:ub-objects} use it: for the entries of the Schur jet matrix of standard data, with
charges $\mathsf w$ and a set $S$ of places at which the columns of height one are removed from the
diagram.

Let $z_\tau$ ($\tau\in\Phi$) be $n'\times n$ matrices of variables $z_{\tau,j,\psi}$ (row $j$, column
$\psi$). A \emph{diagram} $\mathfrak d$ gives the number $\mathfrak d_\tau(h)$ of columns of height
$h$ at $\tau$; only heights $h<n'$ count. A standard filling puts a row index into each box. For
integral $\xi_\tau\in\cO^{n\times n'}$ the \emph{product of minors} $P_\xi$ (\lean{New.minorProd}) is
the product, over the columns, of the minor of $z_\tau\xi_\tau$ on the rows of the column and the
first $h$ columns, $h$ its height. Put $\mathsf a_\tau=\mathfrak d_\tau(1)$,
$\mathsf a'_\tau=\sum_{h\ne1}\mathfrak d_\tau(h)$ and $\mathsf l_\tau=\sum_{h\ne1}h\,\mathfrak d_\tau(h)$
(\lean{New.heightOneColumns}, \lean{New.lowColumnCount}, \lean{New.lowBoxCount}), and let
$B=\sum_\tau(\mathsf a_\tau+\mathsf l_\tau)=\deg P_\xi$ (\lean{New.boxes_eq_sum}) be the number of
boxes (written $L$ in Section~\ref{sec:ub-objects}, Definition~\ref{ub-objects:def:kernel}).

\begin{definition}[Predicted divisibility; \lean{New.predictedDivisibility}]\label{der:def:pred}
For $\chiK(\ell)\ne1$, and with $\tau'=-\ell^e\tau$,
\[
D_\ell(\mathfrak d)=v(\pi_\ell)\sum_{e\in I_\ell}\sum_\tau\Bigl(\bigl(\tfrac{\mathsf a_\tau}{\ell^e}-\mathsf a_{\tau'}
-\mathsf a'_{\tau'}\bigr)_++\bigl(\tfrac{\mathsf l_\tau}{\ell^e}-\mathsf l_{\tau'}\bigr)_+
-\tfrac{\mathsf a'_\tau}{\ell^e}\Bigr)-[\ell=d]\,\tfrac12\sum_\tau\mathrm{ded}(\tau,(-1)\tau,0),
\]
with $\mathrm{ded}(\tau,\tau',e)=\min(\frac{\mathsf a_\tau}{\ell^e},\mathsf a_{\tau'}+\mathsf a'_{\tau'})
+\frac{\mathsf a'_\tau}{\ell^e}+\min(\frac{\mathsf l_\tau}{\ell^e},\mathsf l_{\tau'})$ (\lean{New.deductionTerm}).
\end{definition}

\begin{theorem}[The count; \lean{New.Parameters.norm_jet_minorProd_le_of_hasNormalForm}]
\label{der:thm:count}
Let $\chiK(\ell)\ne1$ and let $F$ have the normal form at $\ell$, with $v(\varpi)=\rho_\ell$, for the
couplings $\Gamma_e=\sum_{\tau,j,\psi}z_{\tau,j,\psi}^{\ell^e}z_{-\ell^e\tau,j,\psi}$. For every
standard filling and integral $\xi$,
$v(\mathrm{jet}_{P_\xi}F)\ge D_\ell(\mathfrak d)-fn'^3(1+\log_\ell\max(B,1))$.
\end{theorem}

\begin{theorem}[Schur entries with charges; \lean{New.Parameters.norm_schurJet_jetSeries_le_of_charges},
\lean{New.Parameters.norm_jet_minorProd_jetSeries_le}]
\label{der:thm:schur}
Let $V=\operatorname{Hom}(E^n,E^{n'})$ be standard, $\mathcal F$ standard data of level one at a
non-split $\ell$, $v(\iota\eta)\ge0$ and $F=\iota J_{\eta\varsigma}\mathcal F$, and let
$\mathfrak d$ fit the ranks ($n\le n'$, no column height strictly between). Let $S\subset\Phi$ and
charges $\mathsf w(\tau,h,j)\in\N$ be such that for $e\in I_\ell^\circ$ and every row $j$: if
$\tau\in S$, or $h=1$ and $-\ell^e\tau\in S$, then $\mathsf w(\tau,1,j)\ne0$ or
$\mathsf w(-\ell^e\tau,h,j)\ne0$. Let $\mathfrak d^\circ$ be $\mathfrak d$ without the columns of
height one at $S$ (\lean{New.eraseHeightOne}), with $B^\circ$ boxes. Then for standard fillings
$\mathfrak t',\mathfrak t$, with $\Sigma\mathsf w$ the sum of the $\mathsf w(\tau,h,j)$ over the boxes
of $\mathfrak t'$ ($h$ the height of the column, $j$ the row),
\[
v\bigl(\mathrm{schurJet}_{\mathfrak d}(F)_{\mathfrak t',\mathfrak t}\bigr)+\Sigma\mathsf w\ \ge\
D_\ell(\mathfrak d^\circ)+\gamma_\ell(B-B^\circ)-fn'^3\bigl(1+\log_\ell\max(B,1)\bigr).
\]
Without charges the same holds for $\mathrm{jet}_{P_\xi}F$.
\end{theorem}

\begin{proof}[Proof of Theorems~\ref{der:thm:count} and~\ref{der:thm:schur}]
By Corollary~\ref{der:cor:integral}(b) the series of Theorem~\ref{der:thm:schur} has the normal form,
and a Schur entry is bounded by the jet along some $P_\xi$, since the $P_\xi$ are combinations of the
Schur entries with coefficients polynomial in $\xi$ and independent modulo the maximal ideal
(\lean{New.exists_norm_schurJet_le_norm_jet_minorProd}). Put
$\mathrm{loss}_B=fn'^3(1+\log_\ell\max(B,1))$, $\mathsf a_\tau:=0$ for $\tau\in S$ (the counts of
$\mathfrak d^\circ$) and
$\Sigma_{\mathrm{ded}}=\sum_{e\in I_\ell^\circ,\,\ell^e\le B}\sum_\tau\mathrm{ded}(\tau,-\ell^e\tau,e)$;
we show $v(\mathrm{jet}_{P_\xi}F)+\Sigma\mathsf w\ge\gamma_\ell B-v(\pi_\ell)\Sigma_{\mathrm{ded}}-\mathrm{loss}_B$.
This suffices: drop the cut ($\mathrm{ded}\ge0$) and split
$\gamma_\ell B=\gamma_\ell(B-B^\circ)+\gamma_\ell B^\circ$; in
$\gamma_\ell B^\circ=v(\pi_\ell)\sum_{e\in I_\ell,\tau}(\mathsf a_\tau+\mathsf l_\tau)\ell^{-e}$ the term
of $(\tau,e)$ less $v(\pi_\ell)\mathrm{ded}(\tau,-\ell^e\tau,e)$ is that of $D_\ell(\mathfrak d^\circ)$,
as $u-\min(u,u')=(u-u')_+$, and $e=0$, in $I_\ell^\circ$ only for $\ell=d$, gives the correction
(\lean{New.mul_tsum_deductionAt}, \lean{New.mul_sum_deductionTerm_le}).

\emph{Step 1: new directions} (\lean{New.exists_newDirection_dvd}). There are $U_{\tau,h}\in M_h(\cO)$
such that the directions $\beta_{\tau hi}=\sum_{i'}(U_{\tau,h})_{ii'}\xi_{\tau,\cdot,i'}$ ($i,i'<h$)
have $\pi_\ell\mid\langle\beta,\beta'\rangle_e=\sum_\psi\beta_\psi^{\ell^e}\beta'_\psi$ for $\beta$ of
height $h$ at $\tau$ and $\beta'$ of height $h'$ at $\tau'=-\ell^e\tau$ ($e\in I_\ell^\circ$,
$\ell^e\le B$), unless $h=1$ and $\tau\in S$, or $h=h'=1$ and $\tau'\in S$, and such that
$\sum_{\tau,h}\mathfrak d_\tau(h)v(\det U_{\tau,h})\le v(\pi_\ell)\Sigma_{\mathrm{ded}}$.
By elimination a linear congruence on all rows of $U_{\tau,h}$ costs the valuation of its modulus once
per column of height $h$. Take it mod $\pi_\ell^{1/\ell^e}$ at $\tau$ ($\ell^e$-th powers are additive
mod $\ell$) and mod $\pi_\ell$ at $\tau'$, on all coordinates of a direction or on its pairing with
the first direction $\xi_{\cdot,0}$ at the other place (a multiple of the one of height one). Covering
the pairs $(1,h')$ costs $\min(\mathsf a_\tau/\ell^e,\mathsf a_{\tau'}+\mathsf a'_{\tau'})$ (coordinates
at $\tau$, or pairings at $\tau'$), the pairs $(h,1)$, $h\ne1$, cost $\mathsf a'_\tau/\ell^e$ (pairings
at $\tau$), and the pairs $(h,h')$, $h,h'\ne1$, cost $\min(\mathsf l_\tau/\ell^e,\mathsf l_{\tau'})$
(coordinates); the sum is $\mathrm{ded}(\tau,\tau',e)$, in units of $v(\pi_\ell)$.

\emph{Step 2: substitution.} Put $z_{\tau,j,\psi}=\sum_{h<n',\,i<h}\ell^{\mathsf w(\tau,h,j)}\beta_{\tau hi\psi}
\tilde z_{\tau hij}$ ($\le fn'^3$ new variables). Modulo $\ell$, the coefficient of
$\tilde z_{\tau hij}^{\ell^e}\tilde z_{\tau'h'i'j}$ in $\Gamma_e(z(\tilde z))$ is
$\ell^{\ell^e\mathsf w(\tau,h,j)+\mathsf w(\tau',h',j)}\langle\beta_{\tau hi},\beta_{\tau'h'i'}\rangle_e$,
divisible by $\pi_\ell$ by Step~1 or, for the exempt pairs, by the charges. So
$G=F(\varpi z(\tilde z))$ is integral in the degrees $\le B$ by Lemma~\ref{der:lem:nf}.

\emph{Step 3: chain rule} (\lean{New.norm_jet_columnDetProd_subst}). For the product $P_*$ over the
columns of $\det(\tilde z_{\tau hij})_{j,i}$ ($j$ the rows of the column, $i<h$), homogeneous of degree
$B$, $\mathrm{jet}_{P_*}G=\varpi^B\ell^{\Sigma\mathsf w}\prod_{\tau,h}(\det U_{\tau,h})^{\mathfrak d_\tau(h)}
\mathrm{jet}_{P_\xi}F$.

\emph{Step 4: Legendre} (\lean{New.norm_jet_le_of_isHomogeneous}). As $G$ is integral in degree $B$,
$v(\mathrm{jet}_{P_*}G)\ge\min_{|\mathbf k|=B}v(\mathbf k!)\ge\frac B{\ell-1}-\mathrm{loss}_B$, each of
the at most $fn'^3$ variables losing at most $1+\log_\ell B$. So
$v(\mathrm{jet}_{P_\xi}F)+\Sigma\mathsf w\ge\frac B{\ell-1}-B\rho_\ell-v(\pi_\ell)\Sigma_{\mathrm{ded}}-\mathrm{loss}_B$,
and $\frac1{\ell-1}-\rho_\ell=\gamma_\ell$.
\end{proof}

\subsection{The predicted divisibility and the prime sums}\label{der:sec:primesums}

This subsection evaluates the predicted divisibility for diagrams of rows and relates it to the
prime sum $\sfin$ of Definition~\ref{setup:def:primesums}: for rows of lengths $\kappa$ and
$\ell\ne d$ it is the local prime sum $S_\ell$, and for $y\le a$, $\sfin(y,\tau)$ is the sum over
the primes $\ell<x_1$ of the local prime sums $s_\ell(y,\tau)$, weighted by $2\log\ell$. This
identifies, up to the error term and the correction at $d$, the divisibility that the count
guarantees with the prime sums of Section~\ref{sec:setup}.

For diagrams of rows ($\mathsf a_\tau$ columns of height one at $\tau$, no others) the predicted
divisibility is
\[
D_\ell=v(\pi_\ell)\sum_\tau\sum_{e\in I_\ell}\bigl(\tfrac{\mathsf a_\tau}{\ell^e}-\mathsf a_{-\ell^e\tau}\bigr)_+
-[\ell=d]\,\tfrac12\sum_\tau\min(\mathsf a_\tau,\mathsf a_{(-1)\tau});
\]
the last term is the term $e=0$ of Definition~\ref{der:def:pred}, which pays for the coupling
$\Gamma_0$ of $\tau$ with $(-1)\tau$.
The \emph{local prime sums} are
$s_\ell(y,\tau)=v(\pi_\ell)\sum_{e\in I_\ell}(y/\ell^e-\kappa(-\ell^e\tau))_+$
(\lean{New.Data.localPrimeSum}; it converges, \lean{New.Data.summable_posPart_sub_weight}) and
$S_\ell=\sum_\tau s_\ell(\kappa(\tau),\tau)$ (\lean{New.Data.localPrimeSumTotal}). If $\mathsf a=\kappa$
at the partners of $\tau$, the term of $\tau$ is $s_\ell(\mathsf a_\tau,\tau)$; so for
$\mathsf a=\kappa$ and $\ell\ne d$, $D_\ell=S_\ell$. For $y\le a$ the prime sum of
Definition~\ref{setup:def:primesums} is
\[
\sfin(y,\tau)=\sum_{\nu\le a}\LambdaK(\nu)\bigl(\tfrac y\nu-\kappa(-\nu\tau)
\bigr)_+=\sum_{\ell<x_1}2\log\ell\cdot s_\ell(y,\tau)
\]
(\lean{New.Data.primeSum_eq_sum_localPrimeSum}): only $\nu=\ell^e$, $e\in I_\ell$, contribute, with
$\LambdaK(\nu)=\Lambda(\nu)(1-\chiK(\nu))=2v(\pi_\ell)\log\ell$. Divisibility by $\ell^{v_\ell}$ saves
$2v_\ell\log\ell$ in the height of a squared norm, so $\sfin(y,\tau)$ is, up to the error term and the
correction at $d$, the saving over all primes that the count guarantees for a jet of order $y$ at
$\tau$ whose partners have the orders $\kappa$. Section~\ref{sec:ub-divis} bounds the correction, for
its diagrams, by $2a/(dm)$ per place, as $\tau$ and $(-1)\tau$ are neither both $a$-places nor both
$b$-places (\lean{New.Data.sub_le_predictedDivisibility}, \lean{New.Data.ne_d_of_placeClass_eq}),
and it compares the diagrams of the composite lift (a row of length $a$ at the $a$-places, columns
of heights $k+i(t+1)$) with the gain (\lean{New.Data.le_divisibilityBound}).

\subsection{The bound is attained}\label{der:sec:sharp}

This subsection shows that the bound of the count is attained up to a normalisation: at every
prime $\ell$, some Taylor coefficient of the jet series of an indicator datum on
$\operatorname{Hom}(E,E^{12})$, multiplied by factorials, by $c_s^{12}$ and by $t_\sharp^{-1}$ for a
number $t_\sharp$ whose conjugates are bounded independently of $\ell$, has valuation at most
$S_\ell$ (Theorem~\ref{der:thm:sharp}).

\begin{theorem}[Sharpness; \lean{New.Data.exists_rpow_neg_localPrimeSumTotal_le_norm_coeff}]
\label{der:thm:sharp}
Let $l\colon\Phi\to\N$ with $\kappa(\tau)\le l_\tau\le\kappa(\tau)+x^{42}$. There is
$t_\sharp\in\C^\times$ with $|\sigma t_\sharp|\le e^{fmb/k}$ for all $\sigma\in\Aut(\C)$ such that for
every modulus $\fm$ with ideal $(\pi_M)$, every prime $\ell$ and every $\iota\mid\ell$ there are an
exponent $\mathbf k$ on $\Phi\times\{1,\dots,12\}$ with $\sum_i\mathbf k(\tau,i)=l_\tau$ and a class
$\mathbf y_0\in(\cO_E/\fm)^{12}$ with
\[
v\bigl(\iota\bigl(c_s^{12}t_\sharp^{-1}\textstyle\prod_\tau l_\tau!\cdot[z^{\mathbf k}]
J_{\pi_M\varsigma}\mathcal F_{\fm,\mathbf y_0}\bigr)\bigr)\le S_\ell,
\]
where $\mathcal F_{\fm,\mathbf y_0}$ on $\operatorname{Hom}(E,E^{12})$ is the vacuum with its component
at the places dividing $\fm$ replaced by the indicator of $\mathbf y_0$ (\lean{New.Data.indicatorDatum}).
\end{theorem}

\emph{Role.} The upper bound does not use this theorem; the lower bound does:
Section~\ref{sec:ol-values} builds from it test forms whose pairing with a lift has valuation at most
$S_\ell+v(c\cdot\mathrm{ratio})$ for a fixed $c\ne0$
(\lean{New.Data.exists_isTestForm_le_norm_inner_trivialLift}).

\begin{proof}[Sketch]
Let $\Theta=\prod_i\theta(\mu_iz)^{c_i}$ along the resolution $w_i/\bar\pi_M$, where $(w_i,c_i)$ is a
tight frame of length $M$ with $c_i\ge6fb$ (\lean{New.Data.exists_isTightFrame}), and let
$\mathrm{val}_n(l)$ be the least valuation of a product of coefficients, of total degree $l$, of $n$
series $H\Theta$, $H$ integral in the $\mathrm y(\mu_iz)$. (1) On one curve $\gamma_\ell$ is optimal:
some $h\in\cO[[\mathrm y]]$ has $v(h^{(n)}(0))\le\gamma_\ell n$, by superadditivity in $n$ and the
formal parameter at $n=q_\ell^j$ (\lean{New.CMLattice.katz_divisibility_limit}). (2) Six copies
$\Theta(\varepsilon_\bullet z)$, $\varepsilon_\bullet\in\Zl$, $\varepsilon_0=1$,
$\sum\varepsilon_\bullet^2=0$, multiply to a unit $U$
(\lean{New.CMLattice.exists_scaledThetaProd_mem_mul_eq_one}); $H_0=U^{-1}h(\mu_{i_0}z)$, where
$\mu_{i_0}$ has unit coefficients, gives $v(l!)+\mathrm{val}_6(l)\le\gamma_\ell|l|$. (3) For inert
$\ell$, in $F(z'+z'')F(\varepsilon_1z')\cdots F(\varepsilon_5z')$, $F=\Theta(\varpi z)=\mathcal W\exp(c_\varpi\Xi)$
(Proposition~\ref{der:prop:thetaprod}), each $z''_{\tau'}$-derivative brings down
$c_\varpi\partial\Xi/\partial z_{\tau'}$, of negative valuation and lowest term
$c_\varpi\lambda_{e0}z_\tau^{\ell^e}$; so $j_\tau$ couplings $z_\tau^{\ell^e}z_{\tau'}$, each lowest
among those with its partner, give $\mathrm{val}_6\le-\sum j_\tau$
(\lean{New.CMLattice.leastValLE_six_of_isCouplingFamily}). Couple $\tau$ with $-\ell\tau$ if $\ell>k$,
and pairs of $b$-places, sparse along the powers of $-\ell$, if $\ell\le k$; six more copies give the
rest of the degree, by (2) or, if $\ell>k$, at valuation $\le0$; for split $\ell$ use (2). So
$v(\prod l_\tau!)+\mathrm{val}_{12}(l)\le S_\ell+2\sum_\tau v(\Pi_\tau)$, where $\Pi_\tau$ is a
product of $\lceil6b/k\rceil$ consecutive integers (\lean{New.Data.leastValLE_twelve}). (4) Twelve
copies of $\prod_i\mathrm y(\mu_iz)^{e_i}\Theta$ form the jet series of a vector with integral values,
and Fourier inversion passes to indicator data (\lean{New.Data.exists_le_norm_coeff_indicatorDatum});
$t_\sharp=M^{12f}c_0\prod_\tau\Pi_\tau^2$, with $c_0\ne0$ a value of $\theta^\circ$.
\end{proof}
